
\documentclass{beamer}
\usetheme{Madrid}
\usecolortheme{default}

\usepackage{amsmath}
\usepackage{amssymb}
\usepackage{bm}
\usepackage{graphicx}
\usepackage{listings}
\usepackage{xcolor}

% Configure Python code formatting
\lstset{
    language=Python,
    basicstyle=\ttfamily\tiny,
    keywordstyle=\color{blue},
    stringstyle=\color{green!50!black},
    commentstyle=\color{red!80!black},
    showstringspaces=false,
    frame=single,
    breaklines=true
}

\title{Constrained Optimization}
\subtitle{Problem 2.0.19 (2007 PI)}
\author{ee25btech11019}
\date{}

\begin{document}

\begin{frame}
    \titlepage
\end{frame}

\begin{frame}{Problem Statement}
    The problem of finding the rectangle of maximum area with perimeter equal to 20 can be posed as the constrained optimization problem:
    \begin{align*}
        \max \quad & xy \\
        \text{s.t.} \quad & 2x + 2y = 20 \\
        & x, y \geq 0
    \end{align*}
    The solution to this problem is $x = y = 5$. What is the value of the Lagrange multiplier corresponding to the perimeter constraint?
    
    \vspace{0.5cm}
    a) 2.5 \hspace{1cm} b) 5 \hspace{1cm} c) 7.5 \hspace{1cm} d) 10
\end{frame}

\begin{frame}{Vector Formulation}
    Let the variables be represented by the vector $\mathbf{x} = \begin{pmatrix} x \\ y \end{pmatrix}$.
    
    The objective function $f(\mathbf{x}) = xy$ can be expressed as a quadratic form:
    \begin{equation}
        f(\mathbf{x}) = \frac{1}{2} \mathbf{x}^\top \mathbf{V} \mathbf{x}
    \end{equation}
    where $\mathbf{V} = \begin{pmatrix} 0 & 1 \\ 1 & 0 \end{pmatrix}$.
    
    \vspace{0.5cm}
    The perimeter constraint $2x + 2y = 20$ can be written as:
    \begin{equation}
        \mathbf{c}^\top \mathbf{x} = d
    \end{equation}
    where $\mathbf{c} = \begin{pmatrix} 2 \\ 2 \end{pmatrix}$ and $d = 20$.
\end{frame}

\begin{frame}{Method of Lagrange Multipliers}
    We formulate the Lagrangian function $L(\mathbf{x}, \lambda)$:
    \begin{equation}
        L(\mathbf{x}, \lambda) = \frac{1}{2} \mathbf{x}^\top \mathbf{V} \mathbf{x} - \lambda (\mathbf{c}^\top \mathbf{x} - d)
    \end{equation}
    where $\lambda$ is the Lagrange multiplier.
    
    \vspace{0.5cm}
    To find the optimum, we set the gradient of $L$ with respect to $\mathbf{x}$ to the zero vector:
    \begin{equation}
        \nabla_{\mathbf{x}} L(\mathbf{x}, \lambda) = \mathbf{V}\mathbf{x} - \lambda\mathbf{c} = \mathbf{0}
    \end{equation}
    \begin{equation}
        \implies \mathbf{x} = \lambda \mathbf{V}^{-1}\mathbf{c}
    \end{equation}
\end{frame}

\begin{frame}{Solving for the Lagrange Multiplier}
    Substitute $\mathbf{x} = \lambda \mathbf{V}^{-1}\mathbf{c}$ into the constraint equation $\mathbf{c}^\top \mathbf{x} = d$:
    \begin{equation}
        \mathbf{c}^\top (\lambda \mathbf{V}^{-1}\mathbf{c}) = d \implies \lambda (\mathbf{c}^\top \mathbf{V}^{-1}\mathbf{c}) = d
    \end{equation}
    
    First, compute the inverse of $\mathbf{V}$:
    \begin{equation}
        \mathbf{V}^{-1} = \begin{pmatrix} 0 & 1 \\ 1 & 0 \end{pmatrix}^{-1} = \begin{pmatrix} 0 & 1 \\ 1 & 0 \end{pmatrix}
    \end{equation}
    
    Next, evaluate the scalar $\mathbf{c}^\top \mathbf{V}^{-1}\mathbf{c}$:
    \begin{equation}
        \mathbf{c}^\top \mathbf{V}^{-1}\mathbf{c} = \begin{pmatrix} 2 & 2 \end{pmatrix} \begin{pmatrix} 0 & 1 \\ 1 & 0 \end{pmatrix} \begin{pmatrix} 2 \\ 2 \end{pmatrix} = \begin{pmatrix} 2 & 2 \end{pmatrix} \begin{pmatrix} 2 \\ 2 \end{pmatrix} = 4 + 4 = 8
    \end{equation}
\end{frame}

\begin{frame}{Final Answer}
    Now, substitute the evaluated scalar back into the constraint equation:
    \begin{equation}
        \lambda (8) = 20
    \end{equation}
    \begin{equation}
        \lambda = \frac{20}{8} = 2.5
    \end{equation}
    
    \vspace{1cm}
    The value of the Lagrange multiplier is 2.5. 
    
    \vspace{0.5cm}
    \textbf{Correct Option: (a)}
\end{frame}

\begin{frame}{Graphical Solution (Simulation)}
    \begin{figure}[H]
        \centering
        % Ensure you run the Python script first to generate this image
        \includegraphics[width=0.8\textwidth]{image.png}
        \caption{Contour plot of Area ($xy$) and Perimeter Constraint ($x+y=10$)}
    \end{figure}
\end{frame}

\begin{frame}[fragile]{Python Code for Simulation}
\begin{lstlisting}
import numpy as np
import matplotlib.pyplot as plt

# Define the objective function (Area)
def f(x, y):
    return x * y

# Generate grid for contour plot
x = np.linspace(0, 10, 400)
y = np.linspace(0, 10, 400)
X, Y = np.meshgrid(x, y)
Z = f(X, Y)

# Plot contours of the objective function
plt.figure(figsize=(8, 6))
contour = plt.contour(X, Y, Z, levels=[5, 10, 15, 20, 25], colors='blue', alpha=0.5)
plt.clabel(contour, inline=True, fontsize=10)

# Plot the perimeter constraint line: 2x + 2y = 20 => y = 10 - x
x_const = np.linspace(0, 10, 400)
y_const = 10 - x_const
plt.plot(x_const, y_const, 'r-', label='Constraint: $x + y = 10$', linewidth=2)

# Plot the optimal point
plt.plot(5, 5, 'ko', markersize=8, label='Optimal Point (5, 5)')

# Formatting the plot
plt.xlabel('$x$')
plt.ylabel('$y$')
plt.title('Constrained Optimization: Maximize Area')
plt.legend()
plt.grid(True)

# Save the plot to include in Beamer
plt.savefig('optimization_plot.png')
plt.show()
\end{lstlisting}
\end{frame}

\end{document}
